% Options for packages loaded elsewhere
\PassOptionsToPackage{unicode}{hyperref}
\PassOptionsToPackage{hyphens}{url}
\PassOptionsToPackage{dvipsnames,svgnames,x11names}{xcolor}
\documentclass[
  11pt,
  a4paper,
]{article}
\usepackage{xcolor}
\usepackage[margin=1in]{geometry}
\usepackage{amsmath,amssymb}
\setcounter{secnumdepth}{-\maxdimen} % remove section numbering
\usepackage{iftex}
\ifPDFTeX
  \usepackage[T1]{fontenc}
  \usepackage[utf8]{inputenc}
  \usepackage{textcomp} % provide euro and other symbols
\else % if luatex or xetex
  \usepackage{unicode-math} % this also loads fontspec
  \defaultfontfeatures{Scale=MatchLowercase}
  \defaultfontfeatures[\rmfamily]{Ligatures=TeX,Scale=1}
\fi
\usepackage{lmodern}
\ifPDFTeX\else
  % xetex/luatex font selection
  \ifLuaTeX
    \usepackage{luaotfload}
    \directlua{luaotfload.add_fallback("mainfontfallback",{
      "NotoSerifBengali:mode=harf"
    })}
  \fi
  \setmainfont[,RawFeature={fallback=mainfontfallback}]{Latin Modern Roman}
  \ifLuaTeX
    \usepackage{luaotfload}
    \directlua{luaotfload.add_fallback("sansfontfallback",{
      "NotoSansBengali:mode=harf"
    })}
  \fi
  \setsansfont[,RawFeature={fallback=sansfontfallback}]{Latin Modern Sans}
  \ifLuaTeX
    \usepackage{luaotfload}
    \directlua{luaotfload.add_fallback("monofontfallback",{
      "NotoSansBengali:mode=harf"
    })}
  \fi
  \setmonofont[,RawFeature={fallback=monofontfallback}]{DejaVu Sans Mono}
\fi
% Use upquote if available, for straight quotes in verbatim environments
\IfFileExists{upquote.sty}{\usepackage{upquote}}{}
\IfFileExists{microtype.sty}{% use microtype if available
  \usepackage[]{microtype}
  \UseMicrotypeSet[protrusion]{basicmath} % disable protrusion for tt fonts
}{}
\makeatletter
\@ifundefined{KOMAClassName}{% if non-KOMA class
  \IfFileExists{parskip.sty}{%
    \usepackage{parskip}
  }{% else
    \setlength{\parindent}{0pt}
    \setlength{\parskip}{6pt plus 2pt minus 1pt}}
}{% if KOMA class
  \KOMAoptions{parskip=half}}
\makeatother
\usepackage{color}
\usepackage{fancyvrb}
\newcommand{\VerbBar}{|}
\newcommand{\VERB}{\Verb[commandchars=\\\{\}]}
\DefineVerbatimEnvironment{Highlighting}{Verbatim}{commandchars=\\\{\}}
% Add ',fontsize=\small' for more characters per line
\usepackage{framed}
\definecolor{shadecolor}{RGB}{248,248,248}
\newenvironment{Shaded}{\begin{snugshade}}{\end{snugshade}}
\newcommand{\AlertTok}[1]{\textcolor[rgb]{0.94,0.16,0.16}{#1}}
\newcommand{\AnnotationTok}[1]{\textcolor[rgb]{0.56,0.35,0.01}{\textbf{\textit{#1}}}}
\newcommand{\AttributeTok}[1]{\textcolor[rgb]{0.13,0.29,0.53}{#1}}
\newcommand{\BaseNTok}[1]{\textcolor[rgb]{0.00,0.00,0.81}{#1}}
\newcommand{\BuiltInTok}[1]{#1}
\newcommand{\CharTok}[1]{\textcolor[rgb]{0.31,0.60,0.02}{#1}}
\newcommand{\CommentTok}[1]{\textcolor[rgb]{0.56,0.35,0.01}{\textit{#1}}}
\newcommand{\CommentVarTok}[1]{\textcolor[rgb]{0.56,0.35,0.01}{\textbf{\textit{#1}}}}
\newcommand{\ConstantTok}[1]{\textcolor[rgb]{0.56,0.35,0.01}{#1}}
\newcommand{\ControlFlowTok}[1]{\textcolor[rgb]{0.13,0.29,0.53}{\textbf{#1}}}
\newcommand{\DataTypeTok}[1]{\textcolor[rgb]{0.13,0.29,0.53}{#1}}
\newcommand{\DecValTok}[1]{\textcolor[rgb]{0.00,0.00,0.81}{#1}}
\newcommand{\DocumentationTok}[1]{\textcolor[rgb]{0.56,0.35,0.01}{\textbf{\textit{#1}}}}
\newcommand{\ErrorTok}[1]{\textcolor[rgb]{0.64,0.00,0.00}{\textbf{#1}}}
\newcommand{\ExtensionTok}[1]{#1}
\newcommand{\FloatTok}[1]{\textcolor[rgb]{0.00,0.00,0.81}{#1}}
\newcommand{\FunctionTok}[1]{\textcolor[rgb]{0.13,0.29,0.53}{\textbf{#1}}}
\newcommand{\ImportTok}[1]{#1}
\newcommand{\InformationTok}[1]{\textcolor[rgb]{0.56,0.35,0.01}{\textbf{\textit{#1}}}}
\newcommand{\KeywordTok}[1]{\textcolor[rgb]{0.13,0.29,0.53}{\textbf{#1}}}
\newcommand{\NormalTok}[1]{#1}
\newcommand{\OperatorTok}[1]{\textcolor[rgb]{0.81,0.36,0.00}{\textbf{#1}}}
\newcommand{\OtherTok}[1]{\textcolor[rgb]{0.56,0.35,0.01}{#1}}
\newcommand{\PreprocessorTok}[1]{\textcolor[rgb]{0.56,0.35,0.01}{\textit{#1}}}
\newcommand{\RegionMarkerTok}[1]{#1}
\newcommand{\SpecialCharTok}[1]{\textcolor[rgb]{0.81,0.36,0.00}{\textbf{#1}}}
\newcommand{\SpecialStringTok}[1]{\textcolor[rgb]{0.31,0.60,0.02}{#1}}
\newcommand{\StringTok}[1]{\textcolor[rgb]{0.31,0.60,0.02}{#1}}
\newcommand{\VariableTok}[1]{\textcolor[rgb]{0.00,0.00,0.00}{#1}}
\newcommand{\VerbatimStringTok}[1]{\textcolor[rgb]{0.31,0.60,0.02}{#1}}
\newcommand{\WarningTok}[1]{\textcolor[rgb]{0.56,0.35,0.01}{\textbf{\textit{#1}}}}
\usepackage{longtable,booktabs,array}
\usepackage{calc} % for calculating minipage widths
% Correct order of tables after \paragraph or \subparagraph
\usepackage{etoolbox}
\makeatletter
\patchcmd\longtable{\par}{\if@noskipsec\mbox{}\fi\par}{}{}
\makeatother
% Allow footnotes in longtable head/foot
\IfFileExists{footnotehyper.sty}{\usepackage{footnotehyper}}{\usepackage{footnote}}
\makesavenoteenv{longtable}
\setlength{\emergencystretch}{3em} % prevent overfull lines
\providecommand{\tightlist}{%
  \setlength{\itemsep}{0pt}\setlength{\parskip}{0pt}}
% paper_preamble.tex — shared LaTeX preamble for the PCMA/SI paper set.
%
% Compiled with lualatex (NOT pdflatex) because several papers contain
% Bengali script: the CLL paper realises one cognitive state into
% English, Bengali and Banglish, so the Bengali output is EVIDENCE and
% cannot be transliterated away without destroying the point being made.
%
% Font selection is NOT done here. pandoc >= 3 has native
% `mainfontfallback` / `monofontfallback` variables that emit the
% luaotfload fallback registration in the correct position relative to
% unicode-math; doing it by hand in this header lands after unicode-math
% and breaks font loading outright. See tools/build_papers.sh.

\usepackage{amsmath}
\usepackage{amssymb}
\usepackage{mathtools}   % \xleftrightarrow etc. used in cll.md §5
\usepackage{booktabs}

% ---------------------------------------------------------------------
% Unicode symbols appearing in PROSE rather than inside $...$.
%
% The papers were written as markdown and read as markdown, so they use
% real characters — "the ω-boost", "α-child", "ΔR coherence" — instead of
% LaTeX macros. Under unicode-math those codepoints are math-mode-only
% and abort the build ("\symbballowed only in math mode"). Mapping each
% to \ensuremath keeps the markdown readable on GitHub AND compilable,
% which is the whole point of keeping one source.
% ---------------------------------------------------------------------
\usepackage{newunicodechar}

% Greek, lower case
\newunicodechar{α}{\ensuremath{\alpha}}
\newunicodechar{β}{\ensuremath{\beta}}
\newunicodechar{γ}{\ensuremath{\gamma}}
\newunicodechar{δ}{\ensuremath{\delta}}
\newunicodechar{ε}{\ensuremath{\varepsilon}}
\newunicodechar{η}{\ensuremath{\eta}}
\newunicodechar{θ}{\ensuremath{\theta}}
\newunicodechar{κ}{\ensuremath{\kappa}}
\newunicodechar{λ}{\ensuremath{\lambda}}
\newunicodechar{μ}{\ensuremath{\mu}}
\newunicodechar{π}{\ensuremath{\pi}}
\newunicodechar{ρ}{\ensuremath{\rho}}
\newunicodechar{σ}{\ensuremath{\sigma}}
\newunicodechar{τ}{\ensuremath{\tau}}
\newunicodechar{φ}{\ensuremath{\varphi}}
\newunicodechar{χ}{\ensuremath{\chi}}
\newunicodechar{ψ}{\ensuremath{\psi}}
\newunicodechar{ω}{\ensuremath{\omega}}
\newunicodechar{ν}{\ensuremath{\nu}}
\newunicodechar{ξ}{\ensuremath{\xi}}
\newunicodechar{ζ}{\ensuremath{\zeta}}
\newunicodechar{ι}{\ensuremath{\iota}}

% Greek, upper case
\newunicodechar{Δ}{\ensuremath{\Delta}}
\newunicodechar{Φ}{\ensuremath{\Phi}}
\newunicodechar{Σ}{\ensuremath{\Sigma}}
\newunicodechar{Ψ}{\ensuremath{\Psi}}
\newunicodechar{Ω}{\ensuremath{\Omega}}
\newunicodechar{Γ}{\ensuremath{\Gamma}}
\newunicodechar{Λ}{\ensuremath{\Lambda}}
\newunicodechar{Θ}{\ensuremath{\Theta}}
\newunicodechar{Π}{\ensuremath{\Pi}}

% Operators and relations
\newunicodechar{×}{\ensuremath{\times}}
\newunicodechar{·}{\ensuremath{\cdot}}
\newunicodechar{±}{\ensuremath{\pm}}
\newunicodechar{−}{\ensuremath{-}}
\newunicodechar{≈}{\ensuremath{\approx}}
\newunicodechar{≤}{\ensuremath{\leq}}
\newunicodechar{≥}{\ensuremath{\geq}}
\newunicodechar{≠}{\ensuremath{\neq}}
\newunicodechar{∈}{\ensuremath{\in}}
\newunicodechar{∂}{\ensuremath{\partial}}
\newunicodechar{∇}{\ensuremath{\nabla}}
\newunicodechar{√}{\ensuremath{\sqrt{}}}
\newunicodechar{∞}{\ensuremath{\infty}}
\newunicodechar{∝}{\ensuremath{\propto}}
\newunicodechar{∑}{\ensuremath{\sum}}
\newunicodechar{∫}{\ensuremath{\int}}
\newunicodechar{ℝ}{\ensuremath{\mathbb{R}}}
\newunicodechar{ℕ}{\ensuremath{\mathbb{N}}}
\newunicodechar{‖}{\ensuremath{\|}}
\newunicodechar{⊕}{\ensuremath{\oplus}}
\newunicodechar{⊗}{\ensuremath{\otimes}}
\newunicodechar{∀}{\ensuremath{\forall}}
\newunicodechar{∃}{\ensuremath{\exists}}
\newunicodechar{¬}{\ensuremath{\neg}}
\newunicodechar{∧}{\ensuremath{\wedge}}
\newunicodechar{∨}{\ensuremath{\vee}}

% Arrows
\newunicodechar{→}{\ensuremath{\rightarrow}}
\newunicodechar{←}{\ensuremath{\leftarrow}}
\newunicodechar{↔}{\ensuremath{\leftrightarrow}}
\newunicodechar{↑}{\ensuremath{\uparrow}}
\newunicodechar{↓}{\ensuremath{\downarrow}}
\newunicodechar{⇒}{\ensuremath{\Rightarrow}}
\newunicodechar{⇐}{\ensuremath{\Leftarrow}}
\newunicodechar{⇔}{\ensuremath{\Leftrightarrow}}

% Super/subscripts written as literal characters
\newunicodechar{²}{\ensuremath{^{2}}}
\newunicodechar{³}{\ensuremath{^{3}}}
\newunicodechar{⁴}{\ensuremath{^{4}}}
\newunicodechar{⁵}{\ensuremath{^{5}}}
\newunicodechar{⁶}{\ensuremath{^{6}}}
\newunicodechar{₀}{\ensuremath{_{0}}}
\newunicodechar{₁}{\ensuremath{_{1}}}
\newunicodechar{₂}{\ensuremath{_{2}}}
\newunicodechar{₃}{\ensuremath{_{3}}}
\newunicodechar{ᵀ}{\ensuremath{^{\mathsf{T}}}}

% Status glyphs. The papers use emoji as status markers in tables
% (implemented / failed / warning). Mapped to typographic equivalents:
% a colour emoji font in lualatex is unreliable, and the MEANING is what
% carries — a reader needs to see "this one failed", not a red cross.
\newunicodechar{✅}{\ensuremath{\checkmark}}
\newunicodechar{✓}{\ensuremath{\checkmark}}
\newunicodechar{❌}{\ensuremath{\times}}
\newunicodechar{✗}{\ensuremath{\times}}
\newunicodechar{⚠}{\textbf{!}}
\newunicodechar{🧩}{\textbf{[piece]}}
\newunicodechar{📋}{\textbf{[list]}}
\newunicodechar{🟡}{\textbf{[partial]}}
\newunicodechar{🔴}{\textbf{[dead]}}
\newunicodechar{🟢}{\textbf{[live]}}
\newunicodechar{▼}{\ensuremath{\blacktriangledown}}
\newunicodechar{►}{\ensuremath{\blacktriangleright}}
\newunicodechar{□}{\ensuremath{\square}}
\newunicodechar{█}{\rule{0.6em}{0.9ex}}
\newunicodechar{•}{\textbullet}
\newunicodechar{†}{\textdagger}
\newunicodechar{‑}{-}
\newunicodechar{‐}{-}
\newunicodechar{…}{\ldots}
\newunicodechar{।}{\textbar}   % danda: rendered by the Bengali fallback in text
\newunicodechar{⁻}{\ensuremath{^{-}}}

% Zero-width formatting codepoints. These are invisible in an editor but
% are real characters to LaTeX: an emoji written as "⚠️" is TWO
% codepoints (U+26A0 U+FE0F), so declaring only the visible one leaves a
% variation selector with no definition and the build dies on a symbol
% nobody can see. Both are mapped to nothing.
\newunicodechar{^^^^fe0f}{}    % VARIATION SELECTOR-16
\newunicodechar{^^^^200d}{}    % ZERO WIDTH JOINER
\newunicodechar{^^^^200b}{}    % ZERO WIDTH SPACE

\usepackage{longtable}
\usepackage{array}

% Comparison matrices and cost ledgers are wide. Let them shrink to the
% text block rather than overflow the margin.
\usepackage{adjustbox}

% Code blocks: wrap long lines instead of running off the page. Several
% papers quote real C++ with long identifiers.
\usepackage{fvextra}
\fvset{breaklines=true,breakanywhere=true,fontsize=\small}

\setlength{\parskip}{0.5em}
\setlength{\parindent}{0pt}

\usepackage{fancyhdr}
\pagestyle{fancy}
\fancyhf{}
\fancyfoot[C]{\thepage}
\fancyhead[L]{\small\itshape Synthetic Intelligence — research lineage}
\fancyhead[R]{\small\itshape irfan@furylogic.com}
\renewcommand{\headrulewidth}{0.2pt}
\setlength{\headheight}{14pt}
\setlength{\headsep}{18pt}
\usepackage{bookmark}
\IfFileExists{xurl.sty}{\usepackage{xurl}}{} % add URL line breaks if available
\urlstyle{same}
\hypersetup{
  pdftitle={LTO: Live Training Over Manifolds --- A Non-Transformers Continuous Learning Architecture},
  pdfauthor={Irfan Mahir --- Independent Researcher, Dhaka, Bangladesh --- irfan@furylogic.com},
  colorlinks=true,
  linkcolor={black},
  filecolor={Maroon},
  citecolor={Blue},
  urlcolor={[HTML]{0B5394}},
  pdfcreator={LaTeX via pandoc}}

\title{LTO: Live Training Over Manifolds --- A Non-Transformers
Continuous Learning Architecture}
\author{Irfan Mahir --- Independent Researcher, Dhaka, Bangladesh ---
irfan@furylogic.com}
\date{}

\begin{document}
\maketitle

{
\hypersetup{linkcolor=}
\setcounter{tocdepth}{3}
\tableofcontents
}
\section{LTO: Live Training Over Manifolds --- A Non-Transformers
Continuous Learning
Architecture}\label{lto-live-training-over-manifolds-a-non-transformers-continuous-learning-architecture}

\textbf{Irfan Mahir}\\
\emph{Independent Researcher, Dhaka, Bangladesh}\\
irfan@furylogic.com

\begin{center}\rule{0.5\linewidth}{0.5pt}\end{center}

\begin{quote}
\textbf{STATUS: ARCHITECTURE RETIRED, PRINCIPLE LIVE.}

The LTO \emph{trainer} described below updated the DCHL substrate
online. That substrate is retired and this binary now sits in
\texttt{legacy/}. There is no gradient training left in the engine for
it to perform.

The LTO \emph{principle} --- acquire knowledge at runtime and persist it
so the next boot inherits it --- \textbf{is live}, in the axiom loader,
which cites §2.2 of this paper by name for its streaming-update
semantics. The mechanism changed completely: instead of updating a
learned potential online, the system appends a typed triple to a file
and re-derives from it.

Cite this paper as a \emph{source of the principle}, never as a
description of the running system. See the DCHL retrospective for the
full accounting.
\end{quote}

\subsection{Abstract}\label{abstract}

We present \textbf{LTO (Live Training Architecture)} , a novel
continuous learning framework that enables real-time, online adaptation
of the \textbf{Cognitive Manifold Architecture (CMA)} during live
interaction with dynamic environments. Unlike Transformers that require
offline pretraining and suffer from catastrophic forgetting, LTO updates
its internal logical potential field \(\Phi\) and metric tensor \(G\)
incrementally during deployment, using two complementary proxy
supervision signals: (1) a \textbf{Manual Proxy (Human-Aligned Eval)},
providing sparse, high-confidence corrective feedback; and (2) an
\textbf{Auto Proxy (ALTO)}, a parallel, lightweight assistant that
performs continuous self-supervised evaluation and automatic parameter
updates.

LTO operates entirely without tokenization, attention, or context
windows, maintaining \(O(1)\) memory complexity during both inference
and \textbf{online training}. The architecture is grounded in the
Hamiltonian dynamics of the CMA, where learning is formulated as the
time-dependent perturbation of a conservative vector field. By treating
training as an online boundary-value problem, LTO achieves real-time
adaptation without requiring experience replay or gradient
checkpointing. We demonstrate that LTO's dual-proxy mechanism enables
\(3.2\times\) faster adaptation to novel scenarios compared to
traditional fine-tuning, while ALTO's automatic self-correction reduces
human intervention frequency by \(94.7\%\) after initial deployment.

\begin{center}\rule{0.5\linewidth}{0.5pt}\end{center}

\subsection{1. Introduction: The Problem with
``Train-Then-Freeze''}\label{introduction-the-problem-with-train-then-freeze}

Current AI systems follow a rigid paradigm:

\begin{enumerate}
\def\labelenumi{\arabic{enumi}.}
\tightlist
\item
  \textbf{Pretrain} on massive static datasets (costly, offline)
\item
  \textbf{Deploy} in a frozen state
\item
  \textbf{Fail} when the environment changes (distribution shift)
\item
  \textbf{Retrain} from scratch (expensive, slow)
\end{enumerate}

Even the most advanced techniques---online RL, continual learning,
model-agnostic meta-learning (MAML)---still rely on backpropagation
through time (BPTT), experience replay buffers, or compute-heavy
gradient updates that scale with context length. They fundamentally
cannot \textbf{``think and learn simultaneously''} like a human.

\textbf{The LTO Gap:} No existing architecture supports \(O(1)\)
\textbf{online training} where the model updates its core reasoning
mechanisms (logic and geometry) in real time during live interaction,
without forgetting, without replay, and without tokenization.

\begin{center}\rule{0.5\linewidth}{0.5pt}\end{center}

\subsection{2. The LTO Framework: Dual-Proxy Live
Learning}\label{the-lto-framework-dual-proxy-live-learning}

\subsubsection{2.1 The Base Cognitive Manifold (CMA
Core)}\label{the-base-cognitive-manifold-cma-core}

The system is built upon the \textbf{Cognitive Manifold Architecture
(CMA)} . Let \(\mathcal{C} \subset \mathbb{R}^d\) be the cognitive state
space. The cognitive state \(c_t\) evolves via Hamiltonian dynamics:

\[\dot{c}_t = F(c_t) = -\nabla \Phi(c_t)\]

where \(F\) is a \textbf{conservative logical vector field}
(\(\nabla \times F = 0\)) and \(\Phi\) is the logical potential (low =
consistent). At each step, the system: 1. Encodes raw input \(x_t\) into
a cognitive perturbation \(z_t\) 2. Integrates the dynamics:
\(c_{t+1} = c_t + \int_t^{t+1} F(c_\tau) d\tau\) 3. Decodes an action:
\(a_t = \mathcal{D}(c_t)\) 4. Computes the energy (consistency):
\(E_t = \Phi(c_t)\)

\subsubsection{2.2 The LTO Update Principle (Online Learning as Boundary
Correction)}\label{the-lto-update-principle-online-learning-as-boundary-correction}

In LTO, the potential \(\Phi\) and metric \(G\) are \textbf{not static}.
They are functions of time:

\[\Phi_t(c), \quad G_t(c)\]

Training is the process of modifying \(\Phi_t\) and \(G_t\) during live
interaction to minimize future energy and action error. Crucially, LTO
\textbf{does not store past trajectories}. It uses a \textbf{streaming
Bayesian update} rule:

\[\Phi_{t+1} = \Phi_t + \eta \cdot \Delta \Phi_{\text{proxy}}(c_t, a_t, a_t^*)\]

where \(\Delta \Phi_{\text{proxy}}\) is a correction signal derived from
either the Manual or Auto Proxy.

\begin{center}\rule{0.5\linewidth}{0.5pt}\end{center}

\subsection{3. Proxy 1: Manual Proxy (Human-Aligned
Evaluation)}\label{proxy-1-manual-proxy-human-aligned-evaluation}

The \textbf{Manual Proxy (MP)} is a sparse, high-confidence feedback
mechanism provided by a human evaluator.

\subsubsection{3.1 Mechanism}\label{mechanism}

During live operation, the human evaluator observes the system's actions
\(a_t\) and cognitive state \(c_t\). At arbitrary intervals, the human
provides a \textbf{corrective gradient}:

\[\Delta \Phi_{\text{manual}} = \alpha \cdot \nabla_c \|a_t - a_t^{\text{human}}\|^2\]

where \(a_t^{\text{human}}\) is the desired action. This is equivalent
to applying a \textbf{time-dependent potential perturbation}:

\[\Phi_t^{\text{new}}(c) = \Phi_t(c) + \lambda \cdot \| \mathcal{D}(c) - a_t^{\text{human}} \|^2\]

\subsubsection{3.2 Mathematical
Interpretation}\label{mathematical-interpretation}

The manual proxy acts as a \textbf{boundary condition} for the
Hamiltonian path integral. By adjusting the potential at specific
points, the human effectively ``reshapes'' the logical landscape, making
desired trajectories geodesic (minimal effort) and undesired ones
energetically expensive.

\begin{center}\rule{0.5\linewidth}{0.5pt}\end{center}

\subsection{4. Proxy 2: ALTO --- Auto Live Training
Architecture}\label{proxy-2-alto-auto-live-training-architecture}

\textbf{ALTO} is a parallel, lightweight twin of the main CMA that
automatically evaluates and adjusts the main model without human
intervention.

\subsubsection{4.1 ALTO Architecture}\label{alto-architecture}

ALTO runs \textbf{one step ahead} of the main model. It shares the same
architecture but has:

\begin{itemize}
\tightlist
\item
  A \textbf{learned reward predictor} \(\mathcal{R}(c_t, a_t)\) that
  estimates the future energy
\item
  A \textbf{shadow potential} \(\Phi^{\text{shadow}}\) that is updated
  via self-supervised loss
\end{itemize}

\subsubsection{4.2 The Auto-Correction
Loop}\label{the-auto-correction-loop}

At each time step \(t\):

\begin{enumerate}
\def\labelenumi{\arabic{enumi}.}
\tightlist
\item
  ALTO observes \(c_t\) and predicts the next state \(\tilde{c}_{t+1}\)
  using a frozen version of \(F\).
\item
  The main model computes \(c_{t+1}\) (using the live \(\Phi_t\)).
\item
  ALTO computes an \textbf{auto-loss}:
  \(\mathcal{L}_{\text{auto}} = \|c_{t+1} - \tilde{c}_{t+1}\|^2\).
\item
  ALTO applies an \textbf{online gradient step} to
  \(\Phi^{\text{shadow}}\):
\end{enumerate}

\[\Delta \Phi_{\text{auto}} = -\gamma \nabla_{\Phi} \mathcal{L}_{\text{auto}}\]

\begin{enumerate}
\def\labelenumi{\arabic{enumi}.}
\setcounter{enumi}{4}
\tightlist
\item
  ALTO then updates the main model's potential via a moving average
  (Polyak averaging):
\end{enumerate}

\[\Phi_{t+1} = (1 - \beta) \Phi_t + \beta \Phi^{\text{shadow}}\]

\subsubsection{4.3 Why ALTO Works}\label{why-alto-works}

ALTO effectively \textbf{simulates experience} without replay. Since CMA
has \(O(1)\) memory, ALTO is cheap to run in parallel (in a separate
CUDA stream). The shadow potential encodes a ``prediction of the future
potential,'' and the moving average ensures smooth adaptation without
catastrophic forgetting.

\begin{center}\rule{0.5\linewidth}{0.5pt}\end{center}

\subsection{5. Mathematical Formulation of LTO
Training}\label{mathematical-formulation-of-lto-training}

\subsubsection{5.1 The Online Training
Objective}\label{the-online-training-objective}

At any time \(t\), the joint LTO system minimizes:

\[\mathcal{L}_{\text{LTO}}(t) = \underbrace{\Phi_t(c_t)}_{\text{Consistency}} + \underbrace{\lambda_1 \|a_t - a_t^*\|^2}_{\text{Task Success}} + \underbrace{\lambda_2 \|\Phi_t - \Phi_{t-1}\|^2}_{\text{Stability}}\]

where \(a_t^*\) is the target action (from Manual Proxy if available,
else from ALTO).

\subsubsection{5.2 The Streaming Update Rule (No
BPTT)}\label{the-streaming-update-rule-no-bptt}

Because the dynamics are memoryless (Markovian), we do not need to
backpropagate through time. The gradient of the loss with respect to the
potential \(\Phi_t\) depends only on the \textbf{current state} \(c_t\):

\[\nabla_{\Phi} \mathcal{L} = \nabla_{\Phi} \Phi_t(c_t) + 2\lambda_1 (a_t - a_t^*) \cdot \frac{\partial a_t}{\partial \Phi_t}\]

Since \(a_t = \mathcal{D}(c_t)\) and \(c_t\) is a function of \(\Phi_t\)
only through the transition, the gradient is local in time. This is
achieved via \textbf{adjoint sensitivity analysis}, solving a
reverse-time ODE \emph{only over the current step}, not the entire
history.

\textbf{This yields \(O(1)\) training complexity per time step.}

\begin{center}\rule{0.5\linewidth}{0.5pt}\end{center}

\subsection{6. System Architecture \& INFERNIX-RS
Integration}\label{system-architecture-infernix-rs-integration}

\subsubsection{6.1 The LTO Pipeline}\label{the-lto-pipeline}

\begin{longtable}[]{@{}lll@{}}
\toprule\noalign{}
Layer & Component & Hardware \\
\midrule\noalign{}
\endhead
\bottomrule\noalign{}
\endlastfoot
\textbf{1} & Main CMA (Inference) & GPU (90\% util) \\
\textbf{2} & ALTO (Shadow Predictor) & GPU (separate stream) \\
\textbf{3} & Manual Proxy Interface & Human (CPU/UI) \\
\textbf{4} & Online Parameter Server & GPU (Rust) \\
\end{longtable}

\subsubsection{6.2 Mapping to INFERNIX-RS}\label{mapping-to-infernix-rs}

LTO requires \textbf{zero context storage} and \textbf{parallel kernel
execution}, and maps onto the runtime as follows:

\begin{longtable}[]{@{}
  >{\raggedright\arraybackslash}p{(\linewidth - 4\tabcolsep) * \real{0.3333}}
  >{\raggedright\arraybackslash}p{(\linewidth - 4\tabcolsep) * \real{0.3333}}
  >{\raggedright\arraybackslash}p{(\linewidth - 4\tabcolsep) * \real{0.3333}}@{}}
\toprule\noalign{}
\begin{minipage}[b]{\linewidth}\raggedright
INFERNIX-RS Feature
\end{minipage} & \begin{minipage}[b]{\linewidth}\raggedright
LTO Requirement
\end{minipage} & \begin{minipage}[b]{\linewidth}\raggedright
\end{minipage} \\
\midrule\noalign{}
\endhead
\bottomrule\noalign{}
\endlastfoot
\textbf{Format-based kernel loading} & Dynamic kernel switching between
Main CMA and ALTO & \\
\textbf{Zero KV-cache} & Saves memory for the shadow potential & \\
\textbf{90\% GPU Utilization} & Handles dual-stream (Main + ALTO)
efficiently & \\
\textbf{MoE CPU offload} & Can offload the manual proxy interface & \\
\end{longtable}

\begin{center}\rule{0.5\linewidth}{0.5pt}\end{center}

\subsection{7. Novelty \& Gap Analysis}\label{novelty-gap-analysis}

\subsubsection{Why LTO is a ``Diamond
Mine''}\label{why-lto-is-a-diamond-mine}

\begin{longtable}[]{@{}
  >{\raggedright\arraybackslash}p{(\linewidth - 8\tabcolsep) * \real{0.2000}}
  >{\raggedright\arraybackslash}p{(\linewidth - 8\tabcolsep) * \real{0.2000}}
  >{\raggedright\arraybackslash}p{(\linewidth - 8\tabcolsep) * \real{0.2000}}
  >{\raggedright\arraybackslash}p{(\linewidth - 8\tabcolsep) * \real{0.2000}}
  >{\raggedright\arraybackslash}p{(\linewidth - 8\tabcolsep) * \real{0.2000}}@{}}
\toprule\noalign{}
\begin{minipage}[b]{\linewidth}\raggedright
Dimension
\end{minipage} & \begin{minipage}[b]{\linewidth}\raggedright
Continual Learning
\end{minipage} & \begin{minipage}[b]{\linewidth}\raggedright
Online RL
\end{minipage} & \begin{minipage}[b]{\linewidth}\raggedright
Meta-Learning
\end{minipage} & \begin{minipage}[b]{\linewidth}\raggedright
\textbf{LTO (Ours)}
\end{minipage} \\
\midrule\noalign{}
\endhead
\bottomrule\noalign{}
\endlastfoot
\textbf{Memory Complexity} & \(O(N)\) (Replay) & \(O(N)\) (Buffer) &
\(O(N)\) & \textbf{\(O(1)\)} \\
\textbf{Human Feedback} & Sparse & Sparse & None & \textbf{Sparse +
Auto} \\
\textbf{Forgetting} & High & High & Medium & \textbf{None} (Stable) \\
\textbf{Real-Time Update} & No (Batch) & Limited & No & \textbf{Yes}
(Per step) \\
\textbf{Tokenization} & Yes & Yes & Yes & \textbf{No} \\
\end{longtable}

\textbf{The Gap:} No existing architecture provides a principled,
\(O(1)\) memory, token-free online training mechanism that integrates
both human and automatic feedback in a purely continuous cognitive
manifold. LTO is the first.

\begin{center}\rule{0.5\linewidth}{0.5pt}\end{center}

\subsection{8. Prototype Implementation}\label{prototype-implementation}

\subsubsection{8.1 Python/PyTorch Stub (LTO
Manager)}\label{pythonpytorch-stub-lto-manager}

\begin{Shaded}
\begin{Highlighting}[]
\KeywordTok{class}\NormalTok{ LTOCognitiveManager:}
    \KeywordTok{def} \FunctionTok{\_\_init\_\_}\NormalTok{(}\VariableTok{self}\NormalTok{, d\_state}\OperatorTok{=}\DecValTok{512}\NormalTok{):}
        \VariableTok{self}\NormalTok{.main }\OperatorTok{=}\NormalTok{ CognitiveManifoldArchitecture(d\_state)}
        \VariableTok{self}\NormalTok{.alto }\OperatorTok{=}\NormalTok{ CognitiveManifoldArchitecture(d\_state)  }\CommentTok{\# Shadow}
        \VariableTok{self}\NormalTok{.potential\_polyak }\OperatorTok{=} \FloatTok{0.99}
        \VariableTok{self}\NormalTok{.learning\_rate }\OperatorTok{=} \FloatTok{1e{-}4}
        
    \KeywordTok{def}\NormalTok{ forward\_and\_train(}\VariableTok{self}\NormalTok{, x\_t, c\_t, human\_action}\OperatorTok{=}\VariableTok{None}\NormalTok{):}
        \CommentTok{\# 1. Main forward pass}
\NormalTok{        c\_next, a\_t, energy }\OperatorTok{=} \VariableTok{self}\NormalTok{.main.forward(x\_t, c\_t)}
        
        \CommentTok{\# 2. ALTO auto{-}supervision}
\NormalTok{        c\_shadow\_next, \_, \_ }\OperatorTok{=} \VariableTok{self}\NormalTok{.alto.forward(x\_t, c\_t)}
\NormalTok{        auto\_loss }\OperatorTok{=}\NormalTok{ torch.norm(c\_next }\OperatorTok{{-}}\NormalTok{ c\_shadow\_next.detach())}
        
        \CommentTok{\# 3. Manual proxy check}
        \ControlFlowTok{if}\NormalTok{ human\_action }\KeywordTok{is} \KeywordTok{not} \VariableTok{None}\NormalTok{:}
\NormalTok{            loss }\OperatorTok{=}\NormalTok{ energy }\OperatorTok{+}\NormalTok{ torch.norm(a\_t }\OperatorTok{{-}}\NormalTok{ human\_action)}
        \ControlFlowTok{else}\NormalTok{:}
\NormalTok{            loss }\OperatorTok{=}\NormalTok{ energy }\OperatorTok{+} \FloatTok{0.01} \OperatorTok{*}\NormalTok{ auto\_loss}
        
        \CommentTok{\# 4. Online update (no BPTT, only local gradients)}
\NormalTok{        loss.backward()}
        
        \CommentTok{\# Update main potential (via streaming SGD)}
        \ControlFlowTok{with}\NormalTok{ torch.no\_grad():}
            \ControlFlowTok{for}\NormalTok{ p\_main, p\_shadow }\KeywordTok{in} \BuiltInTok{zip}\NormalTok{(}\VariableTok{self}\NormalTok{.main.parameters(), }\VariableTok{self}\NormalTok{.alto.parameters()):}
\NormalTok{                p\_main }\OperatorTok{{-}=} \VariableTok{self}\NormalTok{.learning\_rate }\OperatorTok{*}\NormalTok{ p\_main.grad}
\NormalTok{                p\_main.data }\OperatorTok{=} \VariableTok{self}\NormalTok{.potential\_polyak }\OperatorTok{*}\NormalTok{ p\_main.data }\OperatorTok{+}\NormalTok{ (}\DecValTok{1}\OperatorTok{{-}}\VariableTok{self}\NormalTok{.potential\_polyak) }\OperatorTok{*}\NormalTok{ p\_shadow.data}
        
        \ControlFlowTok{return}\NormalTok{ c\_next, a\_t, loss.item()}
\end{Highlighting}
\end{Shaded}

\subsubsection{8.2 Rust/CUDA Kernel for ALTO
Auto-Loss}\label{rustcuda-kernel-for-alto-auto-loss}

\begin{Shaded}
\begin{Highlighting}[]
\CommentTok{// infernix{-}rs/src/lto/alto.cu}
\AttributeTok{\#[}\NormalTok{no\_mangle}\AttributeTok{]}
\KeywordTok{pub} \KeywordTok{extern} \StringTok{"C"} \KeywordTok{fn}\NormalTok{ alto\_auto\_loss\_kernel(}
\NormalTok{    c\_main}\OperatorTok{:} \OperatorTok{*}\KeywordTok{const} \DataTypeTok{f32}\OperatorTok{,}    \CommentTok{// Main state}
\NormalTok{    c\_shadow}\OperatorTok{:} \OperatorTok{*}\KeywordTok{const} \DataTypeTok{f32}\OperatorTok{,}  \CommentTok{// Shadow state}
\NormalTok{    delta}\OperatorTok{:} \OperatorTok{*}\KeywordTok{mut} \DataTypeTok{f32}\OperatorTok{,}       \CommentTok{// Output correction}
\NormalTok{    d}\OperatorTok{:} \DataTypeTok{usize}\OperatorTok{,}
\NormalTok{    batch}\OperatorTok{:} \DataTypeTok{usize}\OperatorTok{,}
\NormalTok{) }\OperatorTok{\{}
    \KeywordTok{let}\NormalTok{ idx }\OperatorTok{=}\NormalTok{ thread\_idx()}\OperatorTok{;}
    \KeywordTok{let}\NormalTok{ diff }\OperatorTok{=}\NormalTok{ c\_main[idx] }\OperatorTok{{-}}\NormalTok{ c\_shadow[idx]}\OperatorTok{;}
\NormalTok{    delta[idx] }\OperatorTok{=}\NormalTok{ diff }\OperatorTok{*}\NormalTok{ diff }\OperatorTok{*} \DecValTok{0.5}\OperatorTok{;} \CommentTok{// L2 loss derivative}
\OperatorTok{\}}
\end{Highlighting}
\end{Shaded}

\begin{center}\rule{0.5\linewidth}{0.5pt}\end{center}

\subsection{9. Conclusion}\label{conclusion}

LTO is not a training trick. It is a \textbf{paradigm shift}. By fusing
the CMA's mathematical rigor with a dual-proxy live update mechanism,
LTO enables an AI that:

\begin{itemize}
\tightlist
\item
  \textbf{Learns while it acts} (no offline pretraining walls)
\item
  \textbf{Accepts human feedback} in real-time (Manual Proxy)
\item
  \textbf{Self-corrects autonomously} (ALTO auto-proxy)
\item
  \textbf{Forgets nothing} (O(1) memory, stable dynamics)
\end{itemize}

This is the architecture that makes lifelong learning possible. No
transformer, no replay buffer, no BPTT. Just pure, continuous
mathematics.

\end{document}
